\begin{thebibliography}{99}
\addcontentsline{toc}{section}{Список литературы}
\small
\bibitem{msmarco} Nguyen T., Rosenberg M., Song X. et al. MS MARCO: A Human Generated MAchine Reading COmprehension Dataset // CoCo Workshop at NIPS. --- 2016. --- arXiv:1611.09268.
\bibitem{nq} Kwiatkowski T., Palomaki J., Redfield O. et al. Natural Questions: A Benchmark for Question Answering Research // Transactions of the ACL. --- 2019. --- Vol.~7. --- P.~452--466.
\bibitem{hotpotqa} Yang Z., Qi P., Zhang S. et al. HotpotQA: A Dataset for Diverse, Explainable Multi-hop Question Answering // Proc. EMNLP. --- 2018. --- P.~2369--2380.
\bibitem{beir} Thakur N., Reimers N., Rücklé A. et al. BEIR: A Heterogeneous Benchmark for Zero-shot Evaluation of Information Retrieval Models // NeurIPS Datasets and Benchmarks Track. --- 2021.
\bibitem{squad} Rajpurkar P., Zhang J., Lopyrev K., Liang P. SQuAD: 100,000+ Questions for Machine Comprehension of Text // Proc. EMNLP. --- 2016. --- P.~2383--2392.
\bibitem{twowiki} Ho X., Duong Nguyen A.-K., Sugawara S., Aizawa A. Constructing A Multi-hop QA Dataset for Comprehensive Evaluation of Reasoning Steps // Proc. COLING. --- 2020.
\bibitem{musique} Trivedi H., Balasubramanian N., Khot T., Sabharwal A. MuSiQue: Multihop Questions via Single-hop Question Composition // Transactions of the ACL. --- 2022. --- Vol.~10.
\bibitem{triviaqa} Joshi M., Choi E., Weld D., Zettlemoyer L. TriviaQA: A Large Scale Distantly Supervised Challenge Dataset for Reading Comprehension // Proc. ACL. --- 2017. --- P.~1601--1611.
\bibitem{lewis2020} Lewis P., Perez E., Piktus A. et al. Retrieval-Augmented Generation for Knowledge-Intensive NLP Tasks // Advances in NeurIPS. --- 2020. --- Vol.~33. --- P.~9459--9474.
\bibitem{drqa} Chen D., Fisch A., Weston J., Bordes A. Reading Wikipedia to Answer Open-Domain Questions // Proc. ACL. --- 2017. --- P.~1870--1879.
\bibitem{gao2023} Gao Y., Xiong Y., Gao X. et al. Retrieval-Augmented Generation for Large Language Models: A Survey. --- 2023. --- arXiv:2312.10997.
\bibitem{bm25} Robertson S., Zaragoza H. The Probabilistic Relevance Framework: BM25 and Beyond // Foundations and Trends in Information Retrieval. --- 2009. --- Vol.~3, №~4. --- P.~333--389.
\bibitem{dpr} Karpukhin V., Oguz B., Min S. et al. Dense Passage Retrieval for Open-Domain Question Answering // Proc. EMNLP. --- 2020. --- P.~6769--6781.
\bibitem{e5} Wang L., Yang N., Huang X. et al. Text Embeddings by Weakly-Supervised Contrastive Pre-training. --- 2022. --- arXiv:2212.03533.
\bibitem{rrf} Cormack G.\,V., Clarke C.\,L.\,A., Büttcher S. Reciprocal Rank Fusion Outperforms Condorcet and Individual Rank Learning Methods // Proc. SIGIR. --- 2009. --- P.~758--759.
\bibitem{ragas} Es S., James J., Espinosa-Anke L., Schockaert S. RAGAS: Automated Evaluation of Retrieval Augmented Generation // Proc. EACL: System Demonstrations. --- 2024.
\bibitem{ares} Saad-Falcon J., Khattab O., Potts C., Zaharia M. ARES: An Automated Evaluation Framework for Retrieval-Augmented Generation Systems // Proc. NAACL. --- 2024.
\bibitem{erag} Salemi A., Zamani H. Evaluating Retrieval Quality in Retrieval-Augmented Generation // Proc. SIGIR. --- 2024.
\bibitem{lostmiddle} Liu N.\,F., Lin K., Hewitt J. et al. Lost in the Middle: How Language Models Use Long Contexts // Transactions of the ACL. --- 2024. --- Vol.~12. --- P.~157--173.
\bibitem{powernoise} Cuconasu F., Trappolini G., Siciliano F. et al. The Power of Noise: Redefining Retrieval for RAG Systems // Proc. SIGIR. --- 2024.
\bibitem{ircot} Trivedi H., Balasubramanian N., Khot T., Sabharwal A. Interleaving Retrieval with Chain-of-Thought Reasoning for Knowledge-Intensive Multi-Step Questions // Proc. ACL. --- 2023.
\bibitem{smucker2007} Smucker M.\,D., Allan J., Carterette B. A Comparison of Statistical Significance Tests for Information Retrieval Evaluation // Proc. CIKM. --- 2007. --- P.~623--632.
\bibitem{dror2018} Dror R., Baumer G., Shlomov S., Reichart R. The Hitchhiker's Guide to Testing Statistical Significance in Natural Language Processing // Proc. ACL. --- 2018. --- P.~1383--1392.
\bibitem{friedman} Friedman M. The Use of Ranks to Avoid the Assumption of Normality Implicit in the Analysis of Variance // JASA. --- 1937. --- Vol.~32, №~200. --- P.~675--701.
\bibitem{page} Page E.\,B. Ordered Hypotheses for Multiple Treatments: A Significance Test for Linear Ranks // JASA. --- 1963. --- Vol.~58, №~301. --- P.~216--230.
\bibitem{cochran} Cochran W.\,G. The Comparison of Percentages in Matched Samples // Biometrika. --- 1950. --- Vol.~37, №~3/4. --- P.~256--266.
\bibitem{mcnemar} McNemar Q. Note on the Sampling Error of the Difference Between Correlated Proportions or Percentages // Psychometrika. --- 1947. --- Vol.~12, №~2. --- P.~153--157.
\bibitem{holm} Holm S. A Simple Sequentially Rejective Multiple Test Procedure // Scandinavian Journal of Statistics. --- 1979. --- Vol.~6, №~2. --- P.~65--70.
\bibitem{efron} Efron B., Tibshirani R.\,J. An Introduction to the Bootstrap. --- New York: Chapman \& Hall, 1993. --- 436~p.
\bibitem{gee} Liang K.-Y., Zeger S.\,L. Longitudinal Data Analysis Using Generalized Linear Models // Biometrika. --- 1986. --- Vol.~73, №~1. --- P.~13--22.
\end{thebibliography}
